\section{Introduction}
\label{sec:intro}

URL-based phishing detection is cheap to deploy. A classifier that reads only the address string can run in a mail gateway, a browser extension or a DNS resolver before any page is fetched, and it does not expose the user to the page content \citep{ma2009,khonji2013,sahoo2017}. The literature on such classifiers keeps growing, from hand-crafted lexical features \citep{garera2007,blum2010,verma2015} to character-level convolutional, recurrent and transformer models \citep{bahnsen2017,le2018urlnet,tajaddodianfar2020,aljofey2020,maneriker2021,bozkir2023}. Surveys list dozens of these systems, many reporting accuracy above 0.95 on public benchmarks \citep{zieni2023,safi2023,basit2021}.

These numbers rest on a small set of reusable corpora. Each corpus pairs a phishing feed, usually PhishTank or OpenPhish, with a legitimate source chosen by the corpus authors: popularity lists \citep{wang2019pdrcnn,prasad2024}, search-engine results \citep{sahingoz2019} or crawls \citep{mamun2016}. The two classes are therefore collected by different pipelines, and a pipeline leaves fingerprints on the strings it records. It decides whether a scheme is written, whether \texttt{www.} is kept, whether an ampersand is HTML-escaped and whether a homepage or a deep page is stored. A classifier that reads raw strings can separate the classes by these fingerprints without learning anything about phishing.

Machine learning has names for this failure. Shortcut learning \citep{geirhos2020} and annotation artefacts \citep{gururangan2018,mccoy2019} describe models that exploit a cue correlated with the label in one dataset but not in the deployment population; medical imaging offers well-documented cases \citep{zech2018,degrave2021,lapuschkin2019}. Security research has its own catalogue of pitfalls, including sampling bias, spurious correlations and data snooping \citep{arp2022,arp2024,sommer2010,pendlebury2019,apruzzese2023}. For phishing, \citet{marchal2018} warned that popularity lists contain domain names without paths, \citet{shirazi2018} found dataset-dependent URL features, and \citet{rashid2024} traced cross-dataset failures partly to prefix and path-length differences. What is missing is a measurement for each widely used corpus: how much of its separability comes from formatting, how much a detector actually uses, and how much of the rest comes from registered domains shared between training and test data.

This paper provides that measurement for seven public corpora. We separate three sources of optimistic bias. \emph{Presentation artefacts} are properties of how a URL string was written down, such as an explicit scheme, a leading \texttt{www.}, a trailing slash or an HTML-escaped ampersand. \emph{Sampling artefacts} arise when the classes are sampled at different granularity, for example homepages of popular domains against deep phishing paths. \emph{Domain leakage} arises when URLs from one registered domain fall on both sides of a split, so that the test partly measures memorisation of hosts. The \texttt{https} scheme sits between these categories, since it is a real server property that collectors record inconsistently, and we treat it separately.

The contributions are as follows.
\begin{itemize}
\item A per-corpus audit of the released files: class-wise presentation and structure statistics, exact and canonical duplicates, label conflicts within and across corpora, and registered-domain overlap. It shows, among other things, that one corpus draws its legitimate class from a few hundred registered domains and that another contains thousands of URLs labelled phishing that a second corpus labels legitimate.
\item A presentation probe with a per-flag ablation and path and homepage subsets, which identifies the flag that separates each corpus and tests whether the signal survives when both classes are sampled at the same granularity.
\item A controlled intervention: the same character detector trained on raw strings, canonical strings and host names, and a neural character model trained on raw and canonical strings, under random and registered-domain-grouped splits, with corrected resampled tests over ten repeated splits for the n-gram detector, and controls for label noise and for the scheme type.
\item Cross-corpus and leave-one-corpus-out transfer with test sets purged against the full training corpus by canonical URL and by registered domain, with ROC-AUC as the primary measure, the precision these detectors would reach at realistic phishing prevalence, and an eight-item reporting checklist.
\item A test of whether post-hoc feature attribution (SHAP) reveals these shortcuts, set against the dependence that the interventions measure.
\end{itemize}

The study does not propose a new detector. Its aim is to make the numbers that detectors report interpretable, which matters to anyone who selects a URL filter by its benchmark score. Section~\ref{sec:related} reviews URL detection, corpus construction and dataset-bias research. Section~\ref{sec:method} describes the corpora, feature groups, canonicalisation and protocols. Section~\ref{sec:results} reports the audit, and Section~\ref{sec:discussion} discusses implications, the checklist and threats to validity.
